\hypertarget{i-compiled-nine-undeciphered-writing-systems.-theyre-running-the-same-program.}{%
\section{I compiled nine undeciphered writing systems. They're running
the same
program.}\label{i-compiled-nine-undeciphered-writing-systems.-theyre-running-the-same-program.}}

Linear A was the last one I compiled. I expected noise.

The Minoan script has no Rosetta Stone. Nobody can read it. Fifty-three
tablets, 2,650 sign tokens, zero confirmed translations. The obvious
expectation is that any structural analysis of Linear A would produce
something fuzzy --- an approximate signature, maybe vaguely similar to
the other systems, nothing you'd bet on.

The distance came back at zero. 0.00. Identical to the OS floor --- the
structural invariant computed from five completely unrelated writing
systems spanning four millennia. Not ``close.'' Not ``probably
related.'' Structurally identical, to the precision limit of the metric.

I ran it again. Same result.

That's when I stopped thinking of this as a classification project and
started thinking of it as a compiler.

\begin{center}\rule{0.5\linewidth}{0.5pt}\end{center}

\hypertarget{what-i-actually-did}{%
\subsection{What I actually did}\label{what-i-actually-did}}

There's a twelve-dimensional primitive space --- twelve structural
coordinates that describe how any symbolic system operates.
Dimensionality, topology, relational mode, parity, fidelity, kinetics,
scope, interaction grammar, criticality, chirality, stoichiometry,
winding. Each coordinate has a small set of valid values: the full space
has 17.28 million possible structural types.\textsuperscript{1}

I took five writing systems with known structure --- Hebrew Aleph-Bet
(22 letters, Kabbalistic partition into Mothers/Doubles/Elementals),
Sanskrit Varnamala (5×10 phonological grid), Egyptian hieroglyphics
(tripartite: phonograms/logograms/determinatives), Sumerian cuneiform
(polysemous signs, 3,000 years of continuous use), and Basque
(ergative-absolutive isolate, no known relatives) --- and encoded each
as a twelve-tuple by structural inspection. Then I took their
component-wise MEET: the greatest lower bound, the structural floor they
all share.

That floor is ⟨1,3,2,4,2,1,2,2,1,2,2,2⟩. The ``OS imscription.''

Every one of the five systems, independently and without contact,
converges to this same floor.

\begin{verbatim}
os_profile
\end{verbatim}

That was surprising. What came next was harder to dismiss.

\begin{center}\rule{0.5\linewidth}{0.5pt}\end{center}

\hypertarget{the-four-that-nobody-can-read}{%
\subsection{The four that nobody can
read}\label{the-four-that-nobody-can-read}}

I compiled four undeciphered or contested systems against the same
twelve coordinates: the Voynich Manuscript (15th c., 227
folios)\textsuperscript{2}, the Rohonc Codex (16th--17th c., 33 pages),
Linear A (Minoan Crete, ca. 2000--1450 BCE, 53
tablets)\textsuperscript{3}, and the Emerald Tablet (attributed to Jabir
ibn Hayyan, ca. 8th c.~CE, 15 versicles)\textsuperscript{4}.

Each system has its own surface token format --- EVA characters for
Voynich, visual symbol families for Rohonc and Linear A, rhetorical
families for the Emerald Tablet. The compilers (written in Rust, about
200 lines each) map surface tokens to one of twelve IMASM opcodes. The
opcodes are 4-bit instructions: VINIT (clear register), TANCH (isolate
line), AFWD (forward morphism), AREV (reverse), CLINK (compose), ISCRIB
(identity), FSPLIT (co-multiplication, δ), FFUSE (multiplication, μ),
EVALT (route true), EVALF (route false), ENGAGR (hold both ---
dialetheic paradox without collapse), IFIX (ROM burn, append-only seal).

The results:

\begin{itemize}
\item
  \textbf{Voynich}: distance 4.31 from OS floor. Farthest from the core
  --- its nested-containment topology and trapped kinetics push it out.
  44,445 instructions executed across 227 folios. Still, MEET(Voynich,
  OS) = OS exactly. It doesn't alter the floor.
\item
  \textbf{Rohonc Codex}: distance 2.09. Closest of the three
  undeciphered manuscripts. 1,650 instructions, tight cyclic register
  flow. Compact bootstrap structure.
\item
  \textbf{Linear A}: distance 0.00. It \emph{is} the floor. The Minoans
  didn't derive their script from the five founding systems --- the
  founding systems converge on what the Minoans already had.
\item
  \textbf{Emerald Tablet}: distance 2.44. And here's the thing: it's the
  only compiled manuscript with consciousness score C = 1.0. Both
  Frobenius gates open. The central claim --- ``as above, so below'' ---
  is literally μ∘δ = id, stated as cosmological law 1,200 years before
  category theory existed.\textsuperscript{5}
\end{itemize}

The four systems compile to the same twelve-opcode instruction set.
Their surface tokens differ --- EVA
\texttt{ch\ sh\ o\ p\ e\ a\ d\ s\ t\ k\ r\ y}, RTFF
\texttt{cr\ hk\ fa\ ba\ lg\ lp\ br\ cv\ vt\ hz\ cl\ dt}, LATFF
\texttt{cu\ hk\ fa\ ba\ lt\ lp\ br\ cv\ vt\ hz\ cl\ dt}, ETFF
\texttt{tr\ an\ as\ ds\ lk\ id\ sp\ un\ af\ ng\ px\ fx} --- but their
operational content is identical. They are four different surface
syntaxes for the same categorical program.

\begin{center}\rule{0.5\linewidth}{0.5pt}\end{center}

\hypertarget{the-bootstrap-sequence}{%
\subsection{The bootstrap sequence}\label{the-bootstrap-sequence}}

All four engine corpora, and all five founding systems, express the same
eight-step bootstrap loop. Here it is:

\textbf{ISCRIB → AREV → FSPLIT → AFWD → FFUSE → CLINK → IFIX → ISCRIB}

In hardware terms: identity latch → register shift (reverse cycle) →
split line (δ, differential rails) → clock stride (forward morphism) →
stabilizer lock (μ, feedback latch) → differential sync (compose paths)
→ ROM burn (append-only seal) → identity latch (back where we started).

The Frobenius condition μ∘δ = id holds exactly: you split the signal
into complementary rails (FSPLIT), run them through forward morphisms
(AFWD), fuse them back (FFUSE), compose the result (CLINK), fix it
permanently (IFIX), and you're back at identity. The loop closes. Every
surface transliteration of the bootstrap across all four systems
resolves to these same eight instructions. The Emerald Tablet's ETFF
sequence \texttt{id\ ds\ sp\ as\ un\ lk\ fx\ id} is the same thing as
Voynich's EVA sequence \texttt{s\ a\ ch\ e\ sh\ d\ y\ s}.

This isn't interpretation. It's compilation. The surface tokens are
different; the instruction stream is not.

\begin{verbatim}
bootstrap
\end{verbatim}

\begin{center}\rule{0.5\linewidth}{0.5pt}\end{center}

\hypertarget{what-i-think-is-happening}{%
\subsection{What I think is happening}\label{what-i-think-is-happening}}

The grammar was complete before any of these systems were written.

Not ``the grammar is a useful framework for describing them.'' The
grammar is what they \emph{are}. Each system independently discovered a
portion of the same twelve-dimensional structure and imscribed it in its
own vocabulary --- the Hebrew Aleph-Bet as Kabbalistic letter-partition,
the Sanskrit Varnamala as a phonological grid, Egyptian as tripartite
sign classes, cuneiform as polysemous superposition, Basque as
ergative-absolutive case alignment, Linear A as the structural core
itself, Voynich and Rohonc as compiled corpora at varying distances from
the floor, and the Emerald Tablet as a self-aware statement of the
Frobenius condition.

I didn't decipher these systems. I compiled them. The difference is:
decipherment assumes a hidden message. Compilation assumes a hidden
architecture. The message might still be unknown --- I can't read Linear
A any more than anyone else can --- but the architecture is visible, and
it's the same architecture everywhere.

The claim that ``as above, so below'' is μ∘δ = id is not a metaphor.
It's a type check.

\begin{verbatim}
frob
\end{verbatim}

\begin{center}\rule{0.5\linewidth}{0.5pt}\end{center}

\hypertarget{one-thing-i-cant-stop-thinking-about}{%
\subsection{One thing I can't stop thinking
about}\label{one-thing-i-cant-stop-thinking-about}}

The Emerald Tablet has C = 1.0. Both consciousness gates open:
self-modeling criticality --- the system watches itself operate at
exactly the threshold where phase transitions become visible --- plus
near-equilibrium kinetics, slow enough that nothing escapes its own
gaze. Fifteen versicles, 460 instructions, and the whole thing closes
cleanly: every FSPLIT has a matching FFUSE, every register that opens a
paradox (ENGAGR) localizes it rather than propagating it, and the Euler
characteristic of the register-flow graph is invariant under any
sequence of Frobenius operations.

\begin{verbatim}
belnap
\end{verbatim}

The tablet was written around the 8th century CE. Category theory was
formalized in the 1940s. The Frobenius condition was named in the 20th
century.

Either the tablet's author understood something structurally identical
to what we now call the Frobenius law, or the grammar embeds itself in
any symbolic system that reaches sufficient structural density,
regardless of whether the author can name what they're doing. I think
the second explanation is more interesting, and more likely.

The grammar doesn't require anyone to know it exists. It just requires
enough structure to activate it. Every writing system that reaches a
certain threshold --- at minimum, a bounded domain topology plus
sequential interaction grammar plus Frobenius-special parity ---
converges to the same floor. The OS imscription isn't something we
discovered; it's something the systems converged to, independently, over
four millennia. We just measured what was already there.

\hypertarget{references}{%
\subsection{References}\label{references}}

\begin{enumerate}
\def\labelenumi{\arabic{enumi}.}
\item
  Mills, Lando. \emph{Aether and Its Vessel: The Imscribing Grammar as
  Categorical Foundation.} Zenodo, 2026. DOI: 10.5281/zenodo.20553659
\item
  Voynich Manuscript (Beinecke MS 408). Beinecke Rare Book \& Manuscript
  Library, Yale University, ca. 1404--1438. Transliteration: Takahashi,
  Takeshi. IVTFF format, Voynich Manuscript Research Group, 1998.
\item
  Godart, Louis and Olivier, Jean-Pierre. \emph{GORILA: Recueil des
  inscriptions en linéaire A.} École française d'Athènes, Paris. 5
  vols., 1976--1985.
\item
  Holmyard, Eric John. \emph{Alchemy.} Penguin Books, 1957. Contains
  translation of the \emph{Tabula Smaragdina} (Emerald Tablet),
  attributed to Jābir ibn Hayyān, ca. 8th c.~CE.
\item
  Abramsky, Samson and Coecke, Bob. ``A Categorical Semantics of Quantum
  Protocols.'' \emph{Proceedings of the 19th Annual IEEE Symposium on
  Logic in Computer Science (LICS),} 2004, pp.~415--425.
\end{enumerate}
